\mxunit{premiere-c09}{Fonction exponentielle et modèles continus}\label{C09}
\begin{keypoint}[Le parcours du chapitre]
\textbf{Objectifs.} Utiliser les propriétés algébriques ; dériver ; étudier une croissance ou décroissance ; relier exponentielle et suites géométriques.

\textbf{Prérequis.} Puissances, suites géométriques, dérivation et signe.

\textbf{Repère annuel.} 2–11 février : 8 h\end{keypoint}
\section*{Diagnostic : six questions, 10 minutes}\label{C09-D}
Répondre sans consulter le cours. Indiquer les questions pour lesquelles une aide est nécessaire. Ce diagnostic ne donne pas lieu à une note.
\begin{enumerate}
\item Simplifier $2^3\times2^4$.
\item Exprimer $3^{-2}$.
\item Résoudre $2x+1=5$.
\item Dériver $x^2+2x$.
\item Une décroissance de 20 \% donne quel coefficient ?
\item Pour une suite géométrique $u_0=5,q=2$, calculer $u_n$.
\end{enumerate}
\paragraph{Après le diagnostic.} Reprendre les notions signalées avec le professeur ; refaire une question analogue, puis expliquer la règle utilisée.
\paragraph{Oral et logique.} Expliquer la différence entre taux continu et coefficient multiplicateur discret ; présenter une inégalité prouvée par minimum d’une différence.
\mxnewpage
\section*{Chercher avant de formaliser}\label{C09-RECH}
Comparer les accroissements de $e^x$ et les valeurs de la fonction autour de 0, avec plusieurs pas positifs. Construire trois étapes d’Euler et expliquer la règle choisie.
\paragraph{Organisation.} Recherche individuelle 5 min, confrontation à deux 5 min, mise en commun 5 min. Garder une trace d’un essai, même non abouti. Une conjecture doit ensuite être justifiée.
\section{Définir une fonction nouvelle}\label{C09-U1}
\begin{keypoint}[Définition et résultats]
La fonction exponentielle est l’unique fonction f dérivable sur $\R$ vérifiant $f\prime=f$ et $f(0)=1$. Son existence et son unicité sont admises. On la note $\exp$, puis $e^x$ avec $e=\exp(1)\approx2{,}71828$. Le nombre e est une constante ; x est la variable. L’exponentielle est strictement positive.
\end{keypoint}
\paragraph{Statut et justification.} Existence et unicité admises conformément au programme. La positivité se déduit, après les propriétés algébriques admises ci-dessous, de $\exp (x)=\exp (\dfrac{x}{2})^2$ et de $\exp(x)\exp(-x)=1$.
\begin{example}$\exp(0)=1$, $\exp(1)=e$, $\exp\prime(0)=1$. La tangente en 0 est donc $y=1+x$. La notation $e^x$ étend les propriétés des puissances à tout exposant réel.\end{example}\label{C09-EX1}
\mxinclude{figures/C09-U1}
\begin{pitfall}$e^0=1$, pas 0. La fonction ne prend jamais la valeur 0, même pour des x négatifs.\end{pitfall}
\begin{application}Quelle est l’équation de la tangente à $e^x$ en 1 ?\end{application}\label{C09-AP1}
\section{Propriétés algébriques et comparaisons}\label{C09-U2}
\begin{keypoint}[Définition et résultats]
Pour tous réels x,y, $e^{x+y}=e^xe^y$, $e^{-x}=\dfrac{1}{e^x}$, $e^{x-y}=\dfrac{e^x}{e^y}$. Pour n entier, $(e^x)^n=e^{nx}$. La fonction est strictement croissante sur $\R$, donc $e^u=e^v\iff u=v$ et $e^u<e^v\iff u<v$.
\end{keypoint}
\paragraph{Statut et justification.} La relation d’addition est admise dans le parcours commun ; elle peut être approfondie via l’unicité. La positivité et $\exp\prime=\exp>0$ donnent la stricte croissance. Celle-ci justifie les équivalences sur équations et inéquations.
\begin{example}$e^{2x}e^{1-x}=e^{x+1}$. $e^{3x-2}=e^{x+4}$ équivaut à $3x-2=x+4$, donc x=3. $e^{2x}\leq1=e^0$ équivaut à $x\leq0$.\end{example}\label{C09-EX2}
\begin{pitfall}$e^{x+y}$ est un produit, jamais en général $e^x+e^y$. On ne résout pas $e^x=2$ avec un logarithme non étudié en première.\end{pitfall}
\begin{application}Simplifier $\dfrac{e^{3x}}{e^{x-1}}$ ; résoudre $e^{x+2}>e^{2x}$.\end{application}\label{C09-AP2}
\section{Dériver et étudier}\label{C09-U3}
\begin{keypoint}[Définition et résultats]
Pour a réel, $t\mapsto e^{at}$ est dérivable et sa dérivée vaut $ae^{at}$. Pour $e^{at+b}=e^be^{at}$, la dérivée est $ae^{at+b}$. Les règles de produit et quotient s’appliquent avec les dérivées déjà connues. Le signe de l’exponentielle simplifie souvent l’étude du signe de la dérivée.
\end{keypoint}
\paragraph{Statut et justification.} La formule de dérivation de $e^{at}$ est admise ; le cas avec b en découle en sortant le facteur constant $e^b$. La dérivation composée générale n’est pas nécessaire ici.
\begin{example}Pour $f(x)=(x+1)e^{-x}$, $f\prime(x)=e^{-x}-(x+1)e^{-x}=-xe^{-x}$. Comme $e^{-x}>0$, la dérivée a le signe de $-x$ : maximum global 1 en 0.\end{example}\label{C09-EX3}
\begin{pitfall}Ne pas oublier le facteur a ni le signe négatif dans $e^{-x}$. Le facteur exponentiel est positif, pas nécessairement croissant.\end{pitfall}
\begin{application}Dériver $g(t)=3e^{-2t}$ et étudier son sens de variation.\end{application}\label{C09-AP3}
\section{Du discret au continu}\label{C09-U4}
\begin{keypoint}[Définition et résultats]
Un modèle $Q(t)=Q_0e^{kt}$, avec $Q_0>0$, représente une croissance si $k>0$, une décroissance si $k<0$, une constance si $k=0$. Pour un pas h fixe, $\dfrac{Q(t+h)}{Q(t)}=e^{kh}$. Les valeurs aux instants entiers forment une suite géométrique de raison $e^k$. L’unité de k est l’inverse de l’unité de temps.
\end{keypoint}
\paragraph{Statut et justification.} Le rapport résulte de $e^{k(t+h)}=e^{kt}e^{kh}$. Le sens de variation se déduit de $Q\prime(t)=kQ_0e^{kt}$ et de la positivité du facteur restant.
\begin{example}$Q(t)=100e^{-0{,}2t}$, t en heures : $Q(0)=100$, $\dfrac{Q(t+1)}{Q(t)}=e^{-0{,}2}$. Une simulation peut chercher le premier nombre entier d’heures où Q passe sous 50 ; l’équation exacte n’est pas résolue par logarithme dans ce chapitre.\end{example}\label{C09-EX4}
\mxinclude{figures/C09-U4}
\begin{pitfall}Le taux relatif par unité de temps discrète est $e^k-1$, pas exactement k ; ils sont proches seulement pour k petit.\end{pitfall}
\begin{application}Pour $Q(t)=5e^{0{,}3t}$, exprimer $Q(n)$ et $\dfrac{Q(t+2)}{Q(t)}$.\end{application}\label{C09-AP4}
\section{Construire numériquement et démontrer une inégalité}\label{C09-U5}
\begin{keypoint}[Définition et résultats]
Pour approcher la solution de $y\prime=y$, $y(0)=1$ avec un pas $h$, Euler remplace l’accroissement par $hy$ : $y_{k+1}=(1+h)y_k$. Sur $[0;1]$ avec $h=\dfrac1n$, l’approximation de $e$ est $\left(1+\dfrac1n\right)^n$. Cette construction est expérimentale. L’inégalité $e^x\geq1+x$ s’obtient par étude de fonction.
\end{keypoint}
\paragraph{Statut et justification.} Pour $g(x)=e^x-1-x$, $g\prime(x)=e^x-1$, négative si x<0 et positive si x>0. Le minimum vaut $g(0)=0$, ce qui prouve l’inégalité ; égalité seulement en 0. La convergence d’Euler n’est pas démontrée ici.
\begin{example}Avec n=2, Euler donne $y_0=1,y_1=1{,}5,y_2=2{,}25$. Avec n=10, $(1{,}1)^{10}\approx2{,}59374$. On compare à e ; on ne confond pas valeur approchée et définition de la fonction.\end{example}\label{C09-EX5}
\begin{pitfall}Un graphique qui montre une courbe au-dessus d’une droite ne prouve pas l’inégalité sur tous les réels.\end{pitfall}
\begin{application}Montrer que $e^{2x}\geq1+2x$ pour tout réel x.\end{application}\label{C09-AP5}
\section*{Parcours d’exercices}\label{C09-PARC}
Commencer par les exercices 1 et 2. Choisir ensuite un exercice 3 ou 4 sans aide. Les exercices 5 à 7 sont une réserve de transfert, de recherche ou de programmation : le professeur en désigne un. L’exercice 8 se reprend à distance.
\begingroup
\mxexercisehead{C09-E1}{Transformer}
Pour tout réel $x$, on considère les expressions\[A=e^xe^{2-x},\qquad B=\dfrac{(e^{2x})^3}{e^x},\qquad C=e^{-x}(e^x+1).\]
\begin{enumerate}[label=\arabic*.]
\item Simplifier $A$ et $B$ en justifiant chaque regroupement d’exposants.
\item Développer puis simplifier $C$. Expliquer pourquoi une somme ne se simplifie pas de la même façon qu’un produit.
\end{enumerate}
\mxexercisehead{C09-E2}{Équations et inéquations}
On considère une équation et une inéquation d’inconnue réelle $x$ :\[e^{2x-1}=e^{x+3}\qquad\text{et}\qquad e^{-3x}\geq e^{x-2}.\]
\begin{enumerate}[label=\arabic*.]
\item Résoudre chacune dans $\R$, en citant la propriété de l’exponentielle qui autorise la comparaison des exposants.
\item Lors de la résolution de l’inéquation, expliquer à quelle étape le sens de l’inégalité change. Donner l’ensemble des solutions sous forme d’intervalle.
\end{enumerate}
\mxexercisehead{C09-E3}{Produit et maximum}
La fonction $f$ est définie sur $[0;4]$ par $f(x)=xe^{-x}$.
\begin{enumerate}[label=\arabic*.]
\item Calculer $f\prime(x)$, puis étudier son signe sur l’intervalle.
\item Dresser le tableau de variations de $f$. En déduire son maximum et l’abscisse où il est atteint.
\item Contrôler que les valeurs aux bornes ne donnent pas un maximum plus grand. Conserver les valeurs exactes.
\end{enumerate}
\mxexercisehead{C09-E4}{Une population modélisée}
Les deux parties suivantes sont indépendantes.\par\textbf{Partie A — Une population.} Une population est modélisée par $P(t)=200e^{0{,}1t}$, où $t\geq0$ est le temps écoulé en années. $P(t)$ est un effectif théorique, qui n’est pas nécessairement entier.
\begin{enumerate}[label=\arabic*.]
\item Déterminer l’effectif initial, puis le coefficient multiplicateur entre les instants $t$ et $t+1$. En déduire le taux annuel exact et son pourcentage arrondi au centième.
\item Expliquer pourquoi ce modèle ne peut pas nécessairement décrire une population réelle indéfiniment.
\item \textbf{Partie B — Deux courbes.} Dans un même repère, esquisser les courbes de $t\mapsto e^{\dfrac{t}{2}}$ et $t\mapsto e^{-\dfrac{t}{2}}$. Préciser leurs variations, leur point commun et leur symétrie. Identifier clairement chaque courbe.
\end{enumerate}
\mxexercisehead{C09-E5}{Recherche d’un seuil}
La quantité d’une substance, en milligrammes, est modélisée à l’instant $t\geq0$ par $Q(t)=80e^{-\dfrac{t}{2}}$, où $t$ est exprimé en heures. Les contrôles ont lieu aux instants entiers $n=0,1,2,\ldots$. On cherche le premier contrôle pour lequel la quantité est strictement inférieure à 10 mg.
\begin{enumerate}[label=\arabic*.]
\item Écrire un algorithme de recherche de ce premier entier. Prévoir une borne d’itérations et préciser ce que le programme renvoie si la borne est atteinte.
\item Exécuter l’algorithme, puis vérifier numériquement la quantité au rang trouvé et au rang précédent.
\item Justifier, à l’aide de la décroissance de $Q$, que ce double contrôle établit qu’il s’agit bien du premier rang. Ne pas utiliser de logarithme.
\end{enumerate}
\mxexercisehead{C09-E6}{Une symétrie}
On considère la fonction $f$ définie sur $\R$ par $f(x)=e^x+e^{-x}$.
\begin{enumerate}[label=\arabic*.]
\item Étudier la parité de $f$ et interpréter le résultat pour sa courbe.
\item Démontrer que $f(x)\geq2$ pour tout réel $x$, en choisissant une méthode. Préciser le cas d’égalité.
\item Proposer un contrôle du résultat par une seconde méthode, à partir d’un carré ou d’une étude de variations.
\end{enumerate}
\mxexercisehead{C09-E7}{Unicité sous une hypothèse}
\textbf{Approfondissement.} On suppose que la fonction exponentielle a déjà été construite. Soit $f$ une fonction dérivable sur $\R$ qui vérifie $f\prime(x)=f(x)$ pour tout réel $x$ et $f(0)=1$. On pose $g(x)=f(x)e^{-x}$.
\begin{enumerate}[label=\arabic*.]
\item Calculer $g\prime$ et expliquer pourquoi $g$ est constante sur $\R$.
\item Déterminer cette constante, puis en déduire l’expression de $f$.
\item Cette démonstration établit-elle l’existence ou l’unicité de la fonction cherchée ? Justifier.
\end{enumerate}
\mxexercisehead{C09-E8}{Reprise différée}
On définit la suite $(u_n)$ par $u_n=e^{2n+1}$ pour tout entier naturel $n$.
\begin{enumerate}[label=\arabic*.]
\item Démontrer que cette suite est géométrique, puis préciser son premier terme et sa raison.
\item Calculer $u_0+u_1+u_2$ sous forme exacte. Comparer le calcul direct et l’utilisation d’une formule de somme.
\end{enumerate}
\paragraph{Travail personnel.} Effectuer uniquement la tâche désignée à l’issue de la séance, dans la progression datée : 10–20 min maximum. Les exercices dépendent des notions déjà étudiées ; certaines questions sont donc différées. Reprendre l’exercice 8 une semaine après le chapitre (10 min). Si le temps prévu est dépassé, noter la difficulté et apporter les essais ; ne pas accumuler toute la réserve.
\endgroup
\mxneed{12}\section*{À faire dans le manuel CQFD}
Ces exercices se traitent dans votre manuel. La progression indique la séance ; ils remplacent le travail personnel prévu ce jour-là (10 à 20 minutes au total).
\paragraph{CQFD, exercice 18 p. 187, questions a, b et c.}\label{C09-M1} Simplifier les produits et quotients en citant la propriété des exponentielles employée.
\paragraph{CQFD, exercice 30 p. 187, questions a, b et c.}\label{C09-M2} Résoudre les trois inéquations et écrire les ensembles de solutions sous forme d’intervalles.
\begingroup\small
\mxnewpage\section*{QCM 1 : comprendre et choisir}\label{C09-Q1}
Cinq questions pédagogiques à réponse unique. Choisir une réponse par question ; expliquer ensuite un distracteur. Ce format de chapitre ne reproduit pas le format officiel de l’EAM.
\begin{enumerate}
\item $e^0=$
\begin{enumerate}[label=\Alph*.,itemsep=2pt]
\item $0$
\item $1$
\item $e$
\item $-1$
\end{enumerate}
\item $e^xe^y=$
\begin{enumerate}[label=\Alph*.,itemsep=2pt]
\item $e^{x+y}$
\item $e^{xy}$
\item $e^x+e^y$
\item $e^{x-y}$
\end{enumerate}
\item Dérivée de $e^{-2x}$ :
\begin{enumerate}[label=\Alph*.,itemsep=2pt]
\item $e^{-2x}$
\item $2e^{-2x}$
\item $-e^{-2x}$
\item $-2e^{-2x}$
\end{enumerate}
\item $e^x=1$ équivaut à
\begin{enumerate}[label=\Alph*.,itemsep=2pt]
\item $x=1$
\item $x=e$
\item $x=0$
\item aucune solution
\end{enumerate}
\item Le signe de $e^{-x}$ est
\begin{enumerate}[label=\Alph*.,itemsep=2pt]
\item négatif si x>0
\item toujours strictement positif
\item nul en 0
\item variable
\end{enumerate}
\end{enumerate}
\endgroup
\begingroup\small
\mxneed{18}\section*{QCM 2 : comprendre et choisir}\label{C09-Q2}
Cinq questions pédagogiques à réponse unique. Choisir une réponse par question ; expliquer ensuite un distracteur. Ce format de chapitre ne reproduit pas le format officiel de l’EAM.
\begin{enumerate}
\item $e^{2x}<e^3$ équivaut à
\begin{enumerate}[label=\Alph*.,itemsep=2pt]
\item $x<3$
\item $x>\dfrac{3}{2}$
\item $x<\dfrac{3}{2}$
\item $x<\sqrt3$
\end{enumerate}
\item $u_n=e^{3n}$ a pour raison
\begin{enumerate}[label=\Alph*.,itemsep=2pt]
\item $e^3$
\item $3$
\item $e$
\item $3e$
\end{enumerate}
\item Tangente à $e^x$ en 0 :
\begin{enumerate}[label=\Alph*.,itemsep=2pt]
\item $y=x$
\item $y=e x$
\item $y=1$
\item $y=x+1$
\end{enumerate}
\item Dérivée de $xe^x$ :
\begin{enumerate}[label=\Alph*.,itemsep=2pt]
\item $e^x$
\item $(x+1)e^x$
\item $xe^x$
\item $xe^{x-1}$
\end{enumerate}
\item $e^x\geq1+x$ a égalité
\begin{enumerate}[label=\Alph*.,itemsep=2pt]
\item pour tout x
\item en 1 seulement
\item en 0 seulement
\item jamais
\end{enumerate}
\end{enumerate}
\subsection*{Vrai ou faux ? Justifier.}\label{C09-VF}
\begin{enumerate}
\item $e^{2x}=2e^x$ pour tout x.
\item La fonction $t\mapsto e^{-3t}$ est décroissante.
\item Un taux continu k=0{,}1 correspond exactement à 10 \% de hausse en une unité de temps.
\end{enumerate}
\begin{keypoint}[Bilan personnel]Je conserve les règles des puissances ; je connais la positivité ; je dérive avec le facteur ; je sépare modèle, calcul exact et simulation.

Pour chaque objectif, indiquer : acquis sans aide ; réussi avec aide ; à retravailler. Écrire un exemple et une erreur à éviter.\end{keypoint}
\endgroup
\ifmxteacher\mxinclude{corriges/C09}\fi
